% !TeX root = ../../Book3.tex
% 纲要标题：### F.1 局部化、局部序与标准基

\section{局部化、局部序与标准基}
\label{b3:sec:F01}

在一个点附近，低阶项比高阶项更能反映局部结构。改变单项式次序后，哪些多项式能够当作单位也会改变，因此必须同时确定计算所在的局部化环。

% 纲要细目（与计划纲要同步）：
% * 比较全局、局部及混合序；正式算例准确识别混合局部化的单位与非单位，不把混合序都当作点局部环
% * 以首单项式理想相等定义标准基，证明分式首项相容与有限首单项式生成；存在性与有效求基分开
% * 完整证明有限阶梯给出局部长度；用同一理想比较局部最小生成元数、首单项式生成元数和阶梯长度
% #### F.1.1 全局序、局部序、混合序与相应局部化
% #### F.1.2 首项理想与标准基的定义
% #### F.1.3 标准基、标准单项式与局部最小生成元

\subsection{全局序、局部序、混合序与相应局部化}
\label{b3:subsec:F01:01}

在多项式环中，常把高次项放在前面；在一个点附近，最低次非零项往往更重要。它决定元素在极大理想的哪一层出现，并给出切锥中的齐次关系。局部项序正是为了先处理这些低阶项。不过，换序的同时必须说明工作环和除法方法也发生了什么变化。

\begin{definition}\label{b3:def:F-monomial-orders}
设 \(k\) 为域，\(n\ge1\)，\(S=k[x_1,\ldots,x_n]\)。若单项式上的全序 \(>\) 与乘法相容，则称它为单项式次序。若每个 \(x_i>1\)，则称 \(>\) 为全局序；若每个 \(x_i<1\)，则称 \(>\) 为局部序；若部分变量大于 \(1\)、部分小于 \(1\)，则称 \(>\) 为混合序。对非零多项式 \(f\in S\)，取其有限个非零项中最大的单项式为 \(\operatorname{LM}_{>}(f)\)，相应的带系数项为 \(\operatorname{LT}_{>}(f)\)。
\end{definition}

第8章的“单项式序”默认要求良序；本附录使用“单项式次序”容纳局部与混合情形。若每个变量都大于 \(1\)，乘法相容性与 Dickson 引理仍保证良序，便回到第8章的定义。例如先将总次数较小的单项式规定为较大，同次再用固定字典序比较，就得到局部次数序。在此序中，\(x-y^2\) 的首项是 \(x\)。这里的大小方向与全局次数序相反：局部次数序先取低次项，全局次数序先取高次项；这一比较只针对次数序，不是任意全局序或局部序的定义。

\ProseParagraphBreak
局部序有下降链
\[
1>x>x^2>x^3>\cdots.
\]
沿此链首单项式越来越小，总次数却越来越大；在局部次数序中，同次单项式之间也可以下降。因此，“首项严格下降”不再保证有限步结束，全局多项式除法所用的良序终止证明在这里失效。

\begin{definition}\label{b3:def:F-order-localization}
设 \(k\) 为域，\(S=k[x_1,\ldots,x_n]\)，\(>\) 是其单项式次序。对非零 \(u\in S\)，记 \(\operatorname{LC}_{>}(u)\) 为首项 \(\operatorname{LT}_{>}(u)\) 的系数。令
\[
U_{>}=\{u\in S:\operatorname{LM}_{>}(u)=1\},
\qquad S_{>}=U_{>}^{-1}S.
\]
称 \(S_{>}\) 为该次序相应的局部化。对 \(f/u\ne0\)，定义
\[
\operatorname{LM}_{>}(f/u)=\operatorname{LM}_{>}(f),
\qquad
\operatorname{LT}_{>}(f/u)
=\operatorname{LT}_{>}(f)/\operatorname{LC}_{>}(u).
\]
\end{definition}
\symidx{\(U_>\)、\(S_>\)：单项式次序确定的乘法集与局部化}

乘法相容性给出非零多项式的公式
\[
\operatorname{LM}_{>}(fg)
=\operatorname{LM}_{>}(f)\operatorname{LM}_{>}(g),
\qquad
\operatorname{LT}_{>}(fg)
=\operatorname{LT}_{>}(f)\operatorname{LT}_{>}(g).
\]
因此 \(U_{>}\) 是乘性集。若 \(f/u=g/v\)，则 \(fv=gu\)；利用 \(\operatorname{LM}_{>}(u)=\operatorname{LM}_{>}(v)=1\) 比较两边首项，得到
\(\operatorname{LM}_{>}(f)=\operatorname{LM}_{>}(g)\) 及
\(\operatorname{LT}_{>}(f)\operatorname{LC}_{>}(v)=\operatorname{LT}_{>}(g)\operatorname{LC}_{>}(u)\)。这说明上述首单项式与首项都不依赖分式表示，乘法公式也随之延伸到 \(S_{>}\)。

\ProseParagraphBreak
全局序时 \(U_{>}=k^\times\)，所以 \(S_{>}=S\)；局部序时 \(U_{>}\) 是常数项非零的多项式，故
\[
S_{>}=k[x_1,\ldots,x_n]_{(x_1,\ldots,x_n)}.
\]
混合序所对应的环要由具体次序计算。例如先按 \(x\) 次数作全局比较，同次再以 \(1>y>y^2>\cdots\) 比较，则 \(U_{>}=k[y]\setminus(y)\)，相应环为 \(k[y]_{(y)}[x]\)。在这个指定的块次序中，\(1-y\) 成为单位，\(1-x\) 仍不可逆，完整核验见\cref{b3:exm:F-mixed-order}；任意混合序的工作环不能一概写成这个环。

\ProseParagraphBreak
本附录按次序选择 \(U_{>}\) 及 \(S_{>}\)，所以从全局序改为局部序或混合序时，还须重新确定工作环及其单位。原理想的生成等式、成员关系和标准基条件，都应在新工作环中解释。

\subsection{首项理想与标准基的定义}
\label{b3:subsec:F01:02}

\begin{definition}\label{b3:def:F-standard-basis}
设 \(k\) 为域，\(S=k[x_1,\ldots,x_n]\)，\(>\) 是其单项式次序，\(R=S_{>}\)，\(I\subset R\) 为理想。其首项理想指多项式环 \(S\) 中的单项式理想
\[
L_{>}(I)=(\operatorname{LM}_{>}(f):0\ne f\in I)\subset S.
\]
若 \(G=(g_1,\ldots,g_s)\) 是 \(I\) 的一个由非零元素组成的有限生成组，并满足
\[
L_{>}(I)=(\operatorname{LM}_{>}(g_1),\ldots,
\operatorname{LM}_{>}(g_s)),
\]
则称 \(G\) 为 \(I\) 关于 \(>\) 的\term{标准基}[standard basis]；零理想取空列作为标准基。本附录采用包含全局 Gröbner 基的这一约定。
\end{definition}
\symidx{\(L_>(I)\)：局部或混合次序下的首项理想}

把首项理想放在 \(S\) 中，可以用同一个指数集合记录被首单项式整除的单项式；原理想及其生成关系则始终在 \(R\) 中理解。由于非零首系数在域 \(k\) 中可逆，用 \(\operatorname{LT}_{>}\) 代替 \(\operatorname{LM}_{>}\) 也生成同一个理想。下面的有限性结论只证明标准基存在，尚未给出从指定生成组计算标准基的算法。

\begin{proposition}\label{b3:prop:F-standard-existence}
设 \(k\) 为域，\(S=k[x_1,\ldots,x_n]\)，\(>\) 为乘法相容的单项式全序，\(R=S_{>}\) 为相应局部化。则每个 \(R\) 的理想都有有限标准基。
\end{proposition}
\begin{proof}
零理想取空列即可。设 \(I\ne0\)。多项式环 \(S\) 是 Noether 环，故单项式理想 \(L_{>}(I)\) 可由有限多个元素的首单项式生成；为这些单项式选择原像 \(g_1,\ldots,g_s\in I\)。环 \(R\) 也是 Noether 环，另取 \(I\) 的一个由非零元素组成的有限生成组并加入该列。所得列生成 \(I\)，其首单项式生成 \(L_{>}(I)\)。

对列中的每个元素写 \(g_i=f_i/u_i\)，其中 \(f_i\in S\)、\(u_i\in U_{>}\)。因为 \(f_i=u_i g_i\) 且 \(u_i\) 在 \(R\) 中可逆，以 \(f_i\) 代替 \(g_i\) 不改变在 \(R\) 中生成的理想。首单项式的乘法公式又给出
\[
\operatorname{LM}_{>}(f_i)
=\operatorname{LM}_{>}(u_i)\operatorname{LM}_{>}(g_i)
=\operatorname{LM}_{>}(g_i).
\]
因此逐个清除分母后得到由多项式组成的有限标准基；清分母可能改变首系数，但不改变首单项式。
\end{proof}

这是存在性证明，没有给出局部序下的终止算法。它既不能保证普通除法结束，也不能从任意生成组直接读出全部首项。例如在 \(k[x,y]_{(x,y)}\) 中，两个生成元的最低次项相消后，可能产生一个此前没有发现的首单项式；\subsecref{b3:subsec:F03:03}将算出这样的例子。

\subsection{标准基、标准单项式与局部最小生成元}
\label{b3:subsec:F01:03}

标准基描述首项，而极小生成组描述 \(I/\mathfrak m I\)。这是两个不同的线性化过程。标准单项式则描述商 \(R/I\) 的元素，不能与前两者混称为“基”而省略所在对象。

\begin{proposition}\label{b3:prop:F-finite-staircase}
设 \(k\) 为域，\(n\ge1\)，\(S=k[x_1,\ldots,x_n]\)，\(R=S_{(x_1,\ldots,x_n)}\)，\(\mathfrak m=(x_1,\ldots,x_n)R\)，取先比较低总次数的局部次数序 \(>\)。设 \(I\subset R\) 满足 \(\mathfrak m^N\subset I\)，其中 \(N\ge1\)，并给出由多项式代表组成的标准基 \(G\)。则不属于 \(L_{>}(I)\) 的单项式在 \(R/I\) 中组成有限的 \(k\)-基。
\end{proposition}
\begin{proof}
所有次数至少为 \(N\) 的单项式都属于 \(L_{>}(I)\)。对任意分式 \(f/u\in R\)，令 \(c=u(0)\ne0\)，写成 \(u=c(1-h)\)，其中 \(h\in(x_1,\ldots,x_n)S\)。有限几何级数恒等式给出
\[
(1-h)(1+h+\cdots+h^{N-1})=1-h^N,
\qquad
u^{-1}\equiv c^{-1}(1+h+\cdots+h^{N-1})
\pmod{\mathfrak m^N}.
\]
于是 \(f/u\) 在模 \(\mathfrak m^N\) 后可用多项式表示，再截去次数至少为 \(N\) 的项，便得到次数小于 \(N\) 的代表。若当前多项式的首单项式属于 \(L_{>}(I)\)，就用 \(G\) 的某个元素消去首项；否则把首项移入余项，并从当前多项式中扣除。每步都丢弃次数至少为 \(N\) 的项。当前首单项式严格下降，而约化只产生同次中更小或次数更高的项；次数小于 \(N\) 的单项式只有有限多个，故过程终止。所得余项仅含标准单项式。由于 \(\mathfrak m^N\subset I\)，每步都保持输入与“当前多项式加已收集余项”在 \(R/I\) 中相等；过程结束时，当前多项式为零，输入与余项之差便属于 \(I\)，证明张成性。

若标准单项式的一个非零线性组合属于 \(I\)，其首单项式也属于 \(L_{>}(I)\)，与每个参与单项式都是标准单项式矛盾。故它们线性无关。
\end{proof}

\begin{example}\label{b3:exm:F-three-bases}
在 \(R=k[x,y]_{(x,y)}\) 中令 \(\mathfrak m=(x,y)R\)，取 \(I=(x^2-y^3,xy)\)，并使用同次按 \(x>y\) 的字典序比较的局部次数序。两个生成元是一个极小生成组，因为其次数 \(2\) 的初始形式 \(x^2,xy\) 线性无关，而 \(\mathfrak m I\) 的元素至少为 \(3\) 阶。因此，\(I\) 所需的最少生成元数为 \(\mu(I)=2\)。可是
\[
y(x^2-y^3)-x(xy)=-y^4,
\]
说明这两个生成元不是局部次数序的标准基。加入 \(y^4\) 后得到标准基，其首单项式为 \(x^2,xy,y^4\)，标准单项式为 \(1,x,y,y^2,y^3\)，具体核验见\cref{b3:exm:F-two-equations}。第三个标准基元素对理想生成是多余的，对首项生成却不可省。三种对象在这个例子中给出不同的数目：
\begin{center}
\begin{tabular}{lll}
\toprule
对象 & 控制的空间或理想 & 本例中的数目 \\
\midrule
极小生成组 & \(I/\mathfrak m I\) 的 \(k\)-基 & \(\mu(I)=2\) \\
局部次数标准基 & \(L_{>}(I)\) 的单项式生成元 & \(3\) \\
标准单项式 & \(R/I\) 的 \(k\)-基 & \(5\) \\
\bottomrule
\end{tabular}
\end{center}
\end{example}

若商不具有有限长度，全部标准单项式通常无限。此时形式除法可能给出一个无限级数余项，而不是有限线性组合。有限长度假设在上面命题中保证了有限的向量空间答案；一般形式情形将在\secref{b3:sec:F03}另行说明。

\begin{example}\label{b3:exm:F-mixed-order}
设 \(k\) 为域，\(S=k[x,y]\)，在单项式上取如下块次序：
\[
x^a y^b>x^{a'}y^{b'}
\quad\Longleftrightarrow\quad
a>a'\ \text{或}\ (a=a'\ \text{且}\ b<b').
\]
这里 \(a,b,a',b'\) 都是非负整数；该次序与乘法相容，且 \(x>1>y\)。任何含 \(x\) 的多项式，其首单项式都含正次幂的 \(x\)，所以不属于 \(U_{>}\)。对非零 \(q(y)\in k[y]\)，首单项式是最低次非零项的单项式，故 \(\operatorname{LM}_{>}(q)=1\) 当且仅当 \(q(0)\ne0\)。因此
\[
U_{>}=k[y]\setminus(y),\qquad
R=U_{>}^{-1}k[x,y]=k[y]_{(y)}[x].
\]
最后一个等式也可直接按分母核验：左边的元素是分母仅含 \(y\)、且在 \(y=0\) 非零的分式；右边的多项式只有有限个系数，对这些系数取公分母便得到同样的分式。

由于 \(\operatorname{LM}_{>}(1-y)=1\)，多项式 \(1-y\) 已在乘法集内，因而在 \(R\) 中可逆。另一方面，\(\operatorname{LM}_{>}(1-x)=x\)，所以 \(1-x\notin U_{>}\)。要证明它在 \(R\) 中仍不可逆，令 \(A=k[y]_{(y)}\)。环 \(A\) 是整环，若存在 \(v\in A[x]\) 满足 \((1-x)v=1\)，则 \(v\ne0\)，而
\[
\deg_x((1-x)v)=1+\deg_x v\ge1,
\]
与右边的 \(x\) 次数为 \(0\) 矛盾。故 \(1-y\in R^\times\)，而 \(1-x\notin R^\times\)。这两个判断使用了指定的块次序及其工作环；若把 \(x\) 也改为局部变量，\(1-x\) 就会成为单位。
\end{example}
